% kernel/engineering_knowledge-DE.tex

\chapter{Wissen zur Bewertung des Kandidaten}
\label{chap:engineering-knowledge}

Vorschlag und Vermessung besitzen nun ausdrückliche Provenienz. Keines von
beiden sagt von selbst, welche Krümmungs-, Überhöhungs-, Geschwindigkeits- oder
Realisierungsbedingungen die Aufgabe bestimmen. Diese Grundlage stammt aus
anwendbarem Engineering Knowledge.

\begin{definition}[Engineering Knowledge]
\label{def:engineering-knowledge}
Engineering Knowledge ist angenommene Ingenieurinformation, die innerhalb
ihrer anwendbaren Qualifikationen zur Interpretation, Bewertung, Beschränkung
oder Entscheidung von Ingenieurfragen verwendet werden darf.
\end{definition}

Im laufenden Fall kann dieses Wissen den angenommenen Vorgängerzustand,
anwendbare Eisenbahnregeln, validierte Berechnungsmethoden, Toleranzen und
autorisierte Betriebsannahmen umfassen. Jedes behält Geltungsbereich,
Gültigkeit und Owner. Eine Regel für einen Geschwindigkeitsbereich oder eine
Realisierung darf nicht nur deshalb übernommen werden, weil ihre Variablen
vertraut heißen.

\begin{definition}[Knowledge Ownership]
\label{def:knowledge-ownership}
Knowledge Ownership ist die Verantwortung, die Engineering Knowledge dem
Konzept oder der Autorität zuordnet, das beziehungsweise die zu Festlegung,
Pflege und Änderung berechtigt ist.
\end{definition}

Ownership verhindert, dass ein Solver-Standardwert zur Projektregel und ein
importierter Wert durch wiederholte Nutzung zur angenommenen Absicht wird.
Beobachtung, Berechnung, Annahme, Vorschlag, Akzeptanz und Entscheidung sind
verschiedene Wahrheitsmodi. Provenienz erklärt, woher eine Aussage stammt;
Annahme und Ownership erklären, ob und in welchem Umfang auf sie vertraut
werden darf.

Die Bewertung kann den Kandidaten nun zum anwendbaren Wissen in Beziehung
setzen. Sie kann die Geometrie als mathematisch konsistent und zugleich für die
Geschwindigkeit ungeeignet, mit Überhöhungsbedingungen unvereinbar, durch zu
alte Evidenz unzureichend gestützt oder einem anderen zulässigen Kandidaten
unterlegen erkennen. Dieser Befund ist Entscheidungsevidenz, nicht die
Entscheidung selbst.

Das letzte Kapitel kehrt zum Übergangsvorschlag zurück und fragt, wer seinen
Ingenieurstatus verändern darf.
